% Observed quantities are inserted by scripts/build_report.py.
\documentclass[11pt,a4paper]{article}
\usepackage[margin=1in]{geometry}
\usepackage{fontspec}
\XeTeXlinebreaklocale "zh"
\XeTeXlinebreakskip=0pt plus 1pt
\setmainfont{SimSun}[BoldFont=SimHei,ItalicFont=SimSun,BoldItalicFont=SimHei]
\newfontfamily\LatinText{Times New Roman}
\newcommand{\LatinWord}[1]{{\LatinText #1}}
\renewcommand{\abstractname}{摘要}
\renewcommand{\refname}{参考文献}
\renewcommand{\tablename}{表}
\renewcommand{\figurename}{图}
\setmonofont{Consolas}[Scale=0.9]
\usepackage{amsmath,mathtools}
\usepackage{unicode-math}
\setmathfont{Cambria Math}
\usepackage[expansion=false]{microtype}
\usepackage{graphicx,booktabs,array,tabularx,longtable,caption,enumitem,placeins}
\usepackage{flafter}
\usepackage[hidelinks]{hyperref}
\setlength{\parindent}{1em}
\setlength{\parskip}{0pt}
\setlength{\emergencystretch}{2em}
\clubpenalty=10000
\widowpenalty=10000
\setlength{\tabcolsep}{4pt}
\renewcommand{\arraystretch}{1.1}
\captionsetup{font=small,labelfont=bf}
\renewcommand{\topfraction}{0.9}
\renewcommand{\bottomfraction}{0.8}
\renewcommand{\textfraction}{0.1}
\renewcommand{\floatpagefraction}{0.75}
\setlist{itemsep=0pt,topsep=3pt}
\newcolumntype{Y}{>{\raggedright\arraybackslash}X}
\newcommand{\ResearchFigure}[3]{\begin{figure}[htbp]\centering
\includegraphics[width=\linewidth]{figures/#1.pdf}\caption{#2}\label{#3}\end{figure}}
\title{鲁棒变量选择的实验评估}
\author{}
\date{2026年10月11日}
\makeatletter
\renewcommand{\maketitle}{\begin{center}{\Large\bfseries\@title\par}\vspace{4pt}
{\small\@date\par}\end{center}\vspace{2pt}}
\makeatother
\begin{document}
\maketitle
\begin{abstract}
本报告考察：通过最小化鲁棒检验损失选择少量变量，能否保留两个总体的差异。四种合成机制分别评价检验功效与信号恢复，独立预测和两个真实数据集提供进一步检验。记录中的鲁棒风险最优并没有一致带来最好的功效或预测表现。另一个小规模搜索对照使完整组合评价次数减少 24.2\%，但逐实例中位时间的总和变化为 +4.3\%。因此，减少内层评价本身不能证明提速。
\end{abstract}

\section{实验问题与共同评价流程}\label{sec:design}
假设十二个变量中只能观测三个，应该选哪三个，才能保留两个总体之间的差异？
我们的方法通过最小化两个经验 Wasserstein 球内的最坏检验损失来选择变量。
实验依次回答：所选变量是否仍能检出分布差异，是否找回造成差异的真实变量，
以及是否有助于预测新的观测。最后再考察求解这个选择问题需要多少计算。

合成统计实验包含三个真实信号变量和九个噪声变量。
第一个总体的各坐标独立服从标准正态分布；第二个总体只改变三个信号坐标，
变化方式见表~\ref{tab:mechanisms}。选择之前随机打乱它们的位置。
每种方法根据每类二十个观测选择三个变量。
变量选择、后续分类器拟合与最终测试使用互不重叠的样本，
选择变量时不能查看最终观测。
每种变化都用新生成的样本，独立重复整个选择与评价流程三十次。
图中的点表示平均结果，区间反映这些重复之间的不确定性；
区间计算方法见第~\ref{sec:parameters} 节。
\begin{table}[htbp]\centering\small
\caption{信号变量的四种差异。其余变量在两个类别中的分布完全相同。}\label{tab:mechanisms}
\begin{tabularx}{\linewidth}{@{}p{29mm}Y@{}}\toprule
差异类型 & 第二个总体的变化\\\midrule
均值 & 每个信号坐标的均值从零变为 $0.9$。\\
方差 & 每个信号坐标的标准差从一变为 $1.8$。\\
相关结构 & 信号坐标之间的两两相关系数为 $0.65$，单个坐标的边缘分布不变。\\
混合分布 & 三个坐标都加上同一个随机正负号乘以 $1.5$，均值仍为零。\\\bottomrule
\end{tabularx}\end{table}
\FloatBarrier

这些实验参考 Wang 等人~\cite{wang} 分别研究计算开销、二样本检验和变量恢复的思路，
但使用本研究自己的数据机制，以及固定的样本量和变量预算。
其中相关结构变化是额外设置的难例：单坐标筛选在总体层面看不到差异。
本报告不声称复现原论文的四种数据模型，
也没有复现其统计实验中样本量、维度和变量预算的扫描。

\subsection{各方法如何选择变量}
六种实际可用的方法都只能选择三个变量。
全局方法与贪心方法优化同一个鲁棒损失；
其他方法分别根据分布差异、预测系数或随机抽取来选变量。
\begin{table}[htbp]\centering\small
\caption{实际可用的变量选择方法。各方法均选择三个变量，实现参数集中于第~\ref{sec:parameters} 节。}\label{tab:selectors}
\begin{tabularx}{\linewidth}{@{}p{32mm}Y@{}}\toprule
方法 & 如何选择变量\\\midrule
全局风险选择 & 搜索变量组合，利用下界排除不可能改进当前结果的组合。\\
鲁棒贪心 & 逐个添加变量，每次选择完整重新拟合后鲁棒损失最小的新增变量。\\
MMD 筛选 & 根据一个核分布差异度量，排序各单坐标的两类分布差异。\\
稀疏加性 MMD & 使用 Wang 等人的线性加性核目标，联合选择稀疏的坐标权重。\\
L1 逻辑回归 & 用带绝对值惩罚的逻辑模型，按系数绝对值排序。\\
随机选择 & 均匀抽取三个变量的组合。\\\bottomrule
\end{tabularx}\end{table}
\FloatBarrier

另外两种参照帮助理解结果。
\emph{已知信号}直接使用三个真实信号变量，这些身份对实际方法不可见。
\emph{全部变量}使用十二个变量，超过了观测预算。
它们分别说明已知信号时的表现，以及保留信号和全部噪声时的表现，
不是预算一致的实际竞争方法。

本研究与 Wang 等人的基线比较只有部分对应：
\begin{table}[htbp]\centering\small
\caption{与 Wang 等人实验的对应范围。未实现的方法不声称已经复现。}\label{tab:baseline-scope}
\begin{tabularx}{\linewidth}{@{}p{45mm}Y@{}}\toprule
原论文的研究内容 & 本报告的评价范围\\\midrule
线性加性核变量选择 & 按论文目标实现；保留学习权重的评价另行报告。\\
二次核与高斯核学习 & 未实现。\\
L0 约束的逻辑回归 & 未实现。本文的 L1 系数排序是额外基线，并非 L0 方法的复现。\\
投影 Wasserstein 变量选择 & 未实现。\\\bottomrule
\end{tabularx}\end{table}
\FloatBarrier

下面的统计结果使用较早搜索实现得到的变量组合。
第~\ref{sec:search} 节在小规模优化实例上单独评价改进后的搜索。
新的计时结果不能证明改进版本的统计表现，或大规模问题上的提速。

\FloatBarrier
\section{所选变量能否检测到分布差异？}\label{sec:statistics}
选择完成后，每个组合都使用同一种二样本检验。
这样比较的是变量本身，而不是同时改变选择方法与后续检验。
检验使用最大均值差异（MMD）：比较总体内部的核相似度与跨总体的核相似度。
这里采用的高斯径向基函数（RBF）核为
$k(x,y)=\exp[-\|x-y\|_2^2/(2b^2)]$。

\emph{功效}是在新的批次中检出真实差异的比例，越大越好。
\emph{假警报率}是在两个批次来自同一总体时仍然拒绝零假设的比例，
不应系统性超过名义上限 $0.05$；较低的保守错误率可以接受。
我们使用无偏的平方 MMD 统计量 $T$ 和置换检验：
\begin{equation}\label{eq:permutation}
p=\frac{1+\#\{T_{\rm perm}\ge T_{\rm observed}-10^{-14}\}}{200},\qquad
\text{若 }p\le0.05\text{ 则拒绝。}
\end{equation}
各方法共用最终观测批次和 $199$ 次标签置换。
核带宽只用选择数据确定，之后固定。
图~\ref{fig:power} 考察选择之后还保留多少可以检出的差异，功效越高越好。两个参照组合帮助区分：问题来自变量选择，还是来自检验本身的难度或噪声变量的干扰。

均值变化时，全局选择的功效为 0.928，单坐标 MMD 筛选为 0.992；配对差为 -6.3 个百分点 [-14.3, -1.3]（$n=30$）。均值变化的全局搜索有 30/30 次达到优化容差，但贪心功效为 0.965。精确求解经验鲁棒目标，并不等于让这个独立检验的功效最大。

\ResearchFigure{power_and_calibration_zh}{共同 RBF 检验的功效，越大越好。前六种方法均选择三个变量。最后两行为参照：已知信号使用真实的三个变量；全部变量使用十二个变量。点为种子均值，误差线为 $95\%$ 自助区间。}{fig:power}
\FloatBarrier
方差与混合变化中，已知信号的功效分别为 0.957 和 1.000，使用全部变量时则为 0.138 和 0.207。保留全部信号并不充分：无关坐标仍可能削弱共同检验。相关结构变化中，即便已知信号的功效也只有 0.218；实际方法选择困难的同时，每类四十个观测的检验本身也不够敏感。

只有假警报率仍受到控制，较高的拒绝比例才有意义。表~\ref{tab:null} 用来自同一总体的两组样本重复共同检验。名义上限为 $0.05$：假警报率不应系统性超过它，较低的保守错误率仍然可以接受。这是检验的假警报率，不是分类器把类别零判错的概率。
\begin{table}[htbp]\centering\small\setlength{\tabcolsep}{3.5pt}
\caption{共同 RBF 检验的假警报率，预定水平为 $0.05$。各训练种子内先对批次取平均，再对三十个种子取平均。已知信号与全部变量为诊断参照。}\label{tab:null}
\begin{tabularx}{\linewidth}{@{}Yrrrr@{}}
\toprule 选择方法 & 均值变化 & 方差变化 & 相关结构变化 & 对称混合 \\
\midrule
全局风险 & 0.052 & 0.040 & 0.048 & 0.055 \\
鲁棒贪心 & 0.053 & 0.047 & 0.060 & 0.055 \\
MMD 筛选 & 0.050 & 0.048 & 0.052 & 0.052 \\
Wang 加性 MMD & 0.050 & 0.050 & 0.047 & 0.058 \\
L1 逻辑回归 & 0.043 & 0.045 & 0.050 & 0.050 \\
随机 & 0.048 & 0.050 & 0.047 & 0.052 \\
已知信号 & 0.053 & 0.050 & 0.055 & 0.048 \\
全部变量 & 0.055 & 0.057 & 0.050 & 0.058 \\
\bottomrule\end{tabularx}\end{table}
\FloatBarrier
\subsection{保留学习到的权重后，检验表现有何变化？}
共同 RBF 评价有意不使用各选择方法自身的权重。Wang 的线性加性方法还会学习带正负号的坐标权重。我们另行评价保留这些学习权重的检验。因此，共同功效评价的是其选出的坐标，原生功效保留有符号加性权重，评价本研究实现的选择与检验组合。任一列都不代表对原论文的完整复现。
\begin{table}[htbp]\centering\small\setlength{\tabcolsep}{3.5pt}
\caption{Wang 所选同一组合的两种检验。共同 RBF 功效只使用坐标；原生功效保留有符号加性权重。数值为种子均值和 $95\%$ 区间，并列出原生假警报率。}\label{tab:native}
\begin{tabularx}{\linewidth}{@{}p{24mm}rrrr@{}}
\toprule 差异 & $n$ & 共同 RBF 功效 & 原生功效 & 原生假警报率 \\
\midrule
均值变化 & 30 & \shortstack{0.982\\{[0.962, 0.997]}} & \shortstack{0.977\\{[0.950, 0.998]}} & \shortstack{0.045\\{[0.025, 0.065]}} \\
方差变化 & 30 & \shortstack{0.328\\{[0.222, 0.435]}} & \shortstack{0.518\\{[0.408, 0.625]}} & \shortstack{0.052\\{[0.033, 0.072]}} \\
相关结构变化 & 30 & \shortstack{0.043\\{[0.025, 0.063]}} & \shortstack{0.045\\{[0.030, 0.060]}} & \shortstack{0.052\\{[0.035, 0.070]}} \\
对称混合 & 30 & \shortstack{0.763\\{[0.643, 0.872]}} & \shortstack{0.945\\{[0.915, 0.970]}} & \shortstack{0.040\\{[0.025, 0.057]}} \\
\bottomrule\end{tabularx}\end{table}
\FloatBarrier
方差与混合变化中，原生检验明显更强。若把共同 RBF 功效当成整个加性方法的表现，会低估它；共同评价的目的，是单独比较所选坐标的作用。


\FloatBarrier
\section{选进多少无关变量、漏掉多少真实信号？}\label{sec:recovery}
检验能发现分布差异，并不意味着找回了每个真实信号变量。
因此还要把所选集合 $S$ 与真实信号集合 $T_*$ 比较。
错误发现比例（FDP）是所选变量中无关变量的占比；
未发现比例（NDP）是真实信号中被遗漏的占比。
二者都越小越好：
\begin{equation}\label{eq:recovery}
\operatorname{FDP}=\frac{|S\setminus T_*|}{\max\{1,|S|\}},\qquad
\operatorname{NDP}=\frac{|T_*\setminus S|}{|T_*|}.
\end{equation}
六种实际方法均有 $|S|=|T_*|=3$，所以每次运行的 FDP 和 NDP 相等。
图中一个数值同时表示二者，不重复画同一组数据。
均匀随机选择的期望 FDP 和 NDP 为 $0.75$，完全恢复时为零。
全部变量参照则有 FDP $0.75$、NDP 零：它保留了三个信号，也保留了九个噪声。
图~\ref{fig:recovery} 接着考察这些坐标是否就是真正的信号变量。FDP/NDP 越小越好；接近随机期望 $0.75$ 时，说明恢复表现与均匀选择三个坐标相差不大。
均值变化中，全局选择的 FDP/NDP 为 0.467，MMD 筛选为 0.244。对称混合中，全局选择的 FDP/NDP 为 0.367，MMD 筛选为 0.289。因此，较高的检出概率仍可能伴随选入无关变量；检验与恢复回答的是不同问题。
\ResearchFigure{recovery_and_prediction_zh}{变量恢复错误，越小越好。六种三变量选择方法的 FDP 与 NDP 相等。参考线为完全恢复的零值和随机期望错误 $0.75$；实际随机方法另画为一个点。点和区间以训练种子为单位。}{fig:recovery}
\FloatBarrier


\FloatBarrier
\section{所选变量能否帮助预测新的观测？}\label{sec:prediction}
接下来固定预测模型，只改变它使用的变量。
每个组合都用独立样本拟合一个高斯核支持向量机（SVM）。
其平衡错误率是两类错误率的平均值，越小越好，$0.5$ 对应随机判断的水平。
图~\ref{fig:prediction} 比较六种实际组合上独立拟合的同一种 SVM。平衡错误率越低，预测越好；这里没有直接沿用鲁棒选择阶段拟合出的检验。
均值变化中，全局选择与 MMD 筛选的错误率分别为 0.315 和 0.279。对称混合中，全局选择与 MMD 筛选的错误率分别为 0.314 和 0.289。鲁棒选择目标并没有一致地为这个不同的预测模型提供最好的输入变量。
\ResearchFigure{prediction_error_zh}{独立 SVM 的平衡错误率，越小越好，$0.5$ 为随机判断水平。六种选择方法使用相同的拟合样本和模型参数。点为种子均值，误差线为 $95\%$ 区间。}{fig:prediction}
\FloatBarrier
\subsection{拟合的鲁棒上界与新样本错误率有何不同}
每个组合还用独立样本重新拟合一个鲁棒分类器。记 $p_0(x)$ 为将观测判为类别零的概率；随机化判断按这个概率抽取类别。它的平衡错误率与 SVM 确定性判断的平衡错误率分别为
\begin{equation}\label{eq:errors}\begin{aligned}
\operatorname{BER}_{\rm hard}&=\tfrac12(\operatorname{FPR}+\operatorname{FNR}),\\
\operatorname{BER}_{\rm random}&=\tfrac12\{\mathbb E_0[1-p_0(X)]+\mathbb E_1[p_0(Y)]\}.
\end{aligned}\end{equation}
FPR 与 FNR 是两类各自的错误判断比例。与确定性分类不同，只要输出概率不是零或一，随机化判断仍可能出错，二者不能混用。表~\ref{tab:head} 比较鲁棒分类器的拟合上界、新样本随机化错误率，以及向另一类平移后的错误率。平移大小为 $0$、$0.25$ 和 $0.50$；其中最大值是实际测得的压力错误，不是对 Wasserstein 球内所有分布取上确界。
\begin{table}[htbp]\centering\small\setlength{\tabcolsep}{3.5pt}
\caption{全局选择变量上重新拟合的鲁棒分类器。错误均为随机化平衡错误率。拟合界为两类总风险上界的一半。压力值为每个种子三次规定平移的最大错误，末列按配对运行统计超过拟合界的次数。}\label{tab:head}
\begin{tabularx}{\linewidth}{@{}p{23mm}rrrrr@{}}
\toprule 差异 & $n$ & 拟合界 & 新样本错误 & 平移最大值 & 超过拟合界 \\
\midrule
均值变化 & 30 & 0.251 & 0.448 & 0.462 & 30/30 \\
方差变化 & 30 & 0.275 & 0.482 & 0.484 & 30/30 \\
相关结构变化 & 30 & 0.287 & 0.497 & 0.501 & 30/30 \\
对称混合 & 30 & 0.248 & 0.423 & 0.426 & 30/30 \\
\bottomrule\end{tabularx}\end{table}
\FloatBarrier
120/120 次配对运行的新样本错误率超过拟合的经验界。这不能称为证书失效：该界针对拟合的经验分布球，并未证明未知总体包含在球内。较低的经验最坏损失不能直接理解成已经验证的总体错误率保证。
\subsection{改变不确定性半径}
表~\ref{tab:radius} 同时改变选择阶段与重新拟合鲁棒分类器的半径，考察整个过程，而不是固定变量后只改分类器。三个半径预先设置，并非根据最终测试结果选择。
\begin{table}[htbp]\centering\small\setlength{\tabcolsep}{3.5pt}
\caption{完整过程的半径敏感性。数值为种子均值与 95\% 区间；错误率为随机化平衡错误率，压力错误在每个种子内对预设平移取最大值。}\label{tab:radius}
\begin{tabularx}{\linewidth}{@{}p{22mm}rrrrr@{}}
\toprule 备择 & $\theta$ & $n$ & 功效 & 新样本错误 & 平移最大值 \\
\midrule
均值变化 & 0.10 & 10 & \shortstack{0.790\\{[0.595, 0.915]}} & \shortstack{0.494\\{[0.491, 0.497]}} & \shortstack{0.496\\{[0.493, 0.499]}} \\
均值变化 & 0.25 & 10 & \shortstack{0.855\\{[0.650, 0.980]}} & \shortstack{0.466\\{[0.450, 0.480]}} & \shortstack{0.474\\{[0.461, 0.485]}} \\
均值变化 & 0.50 & 10 & \shortstack{0.955\\{[0.915, 0.990]}} & \shortstack{0.426\\{[0.411, 0.440]}} & \shortstack{0.442\\{[0.429, 0.453]}} \\
对称混合 & 0.10 & 10 & \shortstack{0.495\\{[0.295, 0.710]}} & \shortstack{0.486\\{[0.478, 0.493]}} & \shortstack{0.487\\{[0.480, 0.493]}} \\
对称混合 & 0.25 & 10 & \shortstack{0.630\\{[0.425, 0.845]}} & \shortstack{0.446\\{[0.417, 0.469]}} & \shortstack{0.448\\{[0.419, 0.471]}} \\
对称混合 & 0.50 & 10 & \shortstack{0.970\\{[0.930, 0.995]}} & \shortstack{0.371\\{[0.352, 0.392]}} & \shortstack{0.378\\{[0.357, 0.401]}} \\
\bottomrule\end{tabularx}\end{table}
\FloatBarrier
这两种机制的记录中，较大半径带来更高功效和较低随机化错误率。这是预设比较的结果，并非一般的单调性保证，也没有独立验证据此选择半径的规则。


\FloatBarrier
\section{真实数据上的表现如何？}\label{sec:real}
WDBC 使用三十个可用变量区分良性与恶性乳腺癌观测；
Ionosphere 使用三十四个变量区分正常与异常雷达回波~\cite{wdbc,ionosphere}。
每种方法仍只保留三个变量。与合成实验一样，后续分类器使用互不重叠的拟合与最终观测。
这里没有已知的真实信号集合，因此评价预测表现，不声称测出了变量恢复准确率。
\begin{table}[htbp]\centering\small\setlength{\tabcolsep}{3.5pt}
\caption{真实数据上的独立共同分类器。数值为均值 $\pm$ 重叠分割标准差。鲁棒错误为随机化 BER，SVM 错误为确定性 BER；搜索受限时使用保留的可行组合。}\label{tab:real}
\begin{tabularx}{\linewidth}{@{}p{18mm}Yrr@{}}
\toprule 数据 & 选择器 & \shortstack{鲁棒\\随机化 BER} & SVM BER \\
\midrule
WDBC & 全局风险 & 0.251\(\pm\)0.046 & 0.081\(\pm\)0.017 \\
WDBC & 鲁棒贪心 & 0.220\(\pm\)0.032 & 0.080\(\pm\)0.021 \\
WDBC & MMD 筛选 & 0.179\(\pm\)0.013 & 0.084\(\pm\)0.017 \\
WDBC & Wang 加性 MMD & 0.180\(\pm\)0.013 & 0.082\(\pm\)0.014 \\
WDBC & L1 逻辑回归 & 0.230\(\pm\)0.034 & 0.078\(\pm\)0.020 \\
WDBC & 随机 & 0.310\(\pm\)0.066 & 0.156\(\pm\)0.078 \\
Ionosphere & 全局风险 & 0.300\(\pm\)0.045 & 0.251\(\pm\)0.094 \\
Ionosphere & 鲁棒贪心 & 0.285\(\pm\)0.035 & 0.202\(\pm\)0.081 \\
Ionosphere & MMD 筛选 & 0.284\(\pm\)0.044 & 0.195\(\pm\)0.053 \\
Ionosphere & Wang 加性 MMD & 0.273\(\pm\)0.042 & 0.181\(\pm\)0.058 \\
Ionosphere & L1 逻辑回归 & 0.284\(\pm\)0.034 & 0.198\(\pm\)0.049 \\
Ionosphere & 随机 & 0.330\(\pm\)0.028 & 0.303\(\pm\)0.069 \\
\bottomrule\end{tabularx}\end{table}
\FloatBarrier
全局搜索在 40/40 次分割中达到时间上限，因此不能将这些组合称为已认证最优解。


\FloatBarrier
\section{减少完整组合求解，是否也减少时间？}\label{sec:search}
评价一个变量组合，需要求解它对应的连续鲁棒检验问题。
我们同时统计完整评价了多少个组合，以及选择和评价组合的总时间。
前者反映昂贵的内层工作；后者还包含决定下一个组合时使用的整数搜索。
改进保持原本预测下一个评价组合的规则，只额外保留每次完整求解已经返回的较强对偶分式下界，用它排除更多不值得评价的组合。不增加额外连续求解，任何新选中的二元组合仍然完整评价。对照为原来的高斯过程引导切线搜索，以及每一步都刷新切线主问题的实现。

独立确认研究在 4 份生成数据上设置 16 个配置，维度为七或十，每类六个观测。各方法在每个配置下计时 3 次，下面的总数先在各配置内取中位数再求和。这些设置共享数据，不是 16 次独立的总体试验。表~\ref{tab:search-confirmation} 将完整评价次数与整数主问题的时间分开。
\begin{table}[htbp]\centering\small\setlength{\tabcolsep}{3.5pt}
\caption{当前确认比较。次数与秒数是逐配置中位数的总和，成功运行数包含全部三次计时。每次完整评价执行两个连续求解。}\label{tab:search-confirmation}
\begin{tabularx}{\linewidth}{@{}Yrrrr@{}}
\toprule 搜索 & \shortstack{评价\\组合数} & \shortstack{主问题\\秒数} & \shortstack{总\\秒数} & 最优运行 \\
\midrule
每步切线主问题 & 601 & 13.507 & 17.312 & 48/48 \\
原 GP 搜索 & 598 & 8.271 & 12.715 & 48/48 \\
保留分式 & 453 & 9.610 & 13.262 & 48/48 \\
\bottomrule\end{tabularx}\end{table}
\FloatBarrier
改进使评价次数从 598 降到 453，减少 24.2\%。评价次数减少/相同/增加的配置分别为 15/0/1。总时间从 12.715 变为 13.262 秒（+4.3\%），在 8/16 个配置上更快。整数主问题更昂贵，所以减少的内层评价在这里没有带来整体提速。未改变的 GP 不能被视为这次评价减少的原因。
图~\ref{fig:search-confirmation} 展示配对配置，而不仅给出总数：评价次数越少越好，时间比大于一则表示更慢。
\ResearchFigure{search_confirmation_zh}{完整组合评价次数（上图）与改进/原始时间比（下图）。虚线一表示时间相同。配置 1--8 使用半合页损失，9--16 使用逻辑损失；各组内依次按维度、种子和半径对排列。}{fig:search-confirmation}
\FloatBarrier
确认实验有 144/144 次运行达到数值目标，相对全部组合枚举下界的最大可行风险超出为 \(1.6\times10^{-7}\)。开发中曾有一次主问题被报告为不可行；采用共同的数值修复后，固定规则通过了 26/26 个集成实例。全部开发对照和失败运行保留在 \path{docs/experiment_development_record.md} 中。新的效率证据仅限于这些小规模实例。


\subsection{维度增加时的较早比较}\label{sec:compute}
较早实现还比较了连接二元选择与连续问题的三种价格约束：
限制各坐标的单独价格，限制价格的加权总和，以及用待最小化的风险来缩放该总和上限。
最后一种另行比较使用和不使用分数起始阶段的情况。
四种实现求解相同的二元鲁棒目标。
在 159 对已解出的约束形式/枚举结果中，最大记录可行目标差为 0.000；按合并求解容差核查，0 个区间不相容。在 $D\ge20$ 时，全局风险选择在均值变化的 15 个实例中有 15 个达到时间上限。在 $D\ge20$ 时，全局风险选择在对称混合的 15 个实例中有 15 个达到时间上限。
\begin{table}[htbp]\centering\small\setlength{\tabcolsep}{3.5pt}
\caption{$D=12$ 的参考比较。时间为累计秒数及 95\% 种子自助区间；组合数仅计不同整数子集，不含分数根节点调用。}\label{tab:reference-computation}
\begin{tabularx}{\linewidth}{@{}p{22mm}Yrrr@{}}
\toprule 备择 & 搜索 & 解出数 & 时间（秒） & 组合数 \\
\midrule
均值变化 & 逐项价格界 & 10/10 & \shortstack{4.07\\{[3.57, 4.88]}} & 220.0 \\
均值变化 & 总价格界 & 10/10 & \shortstack{4.63\\{[4.16, 5.08]}} & 219.9 \\
均值变化 & 风险（无根节点） & 10/10 & \shortstack{13.73\\{[13.05, 14.32]}} & 203.2 \\
均值变化 & 全局风险 & 10/10 & \shortstack{16.54\\{[15.74, 17.28]}} & 197.8 \\
对称混合 & 逐项价格界 & 10/10 & \shortstack{3.52\\{[3.40, 3.63]}} & 220.0 \\
对称混合 & 总价格界 & 10/10 & \shortstack{4.75\\{[4.38, 5.08]}} & 219.2 \\
对称混合 & 风险（无根节点） & 10/10 & \shortstack{13.22\\{[12.30, 14.07]}} & 200.0 \\
对称混合 & 全局风险 & 10/10 & \shortstack{16.39\\{[15.24, 17.52]}} & 197.0 \\
\bottomrule\end{tabularx}\end{table}
\FloatBarrier
\ResearchFigure{optimization_comparison_zh}{较早实现：随维度增加的实际时间与解出比例。时间越低、解出比例越高越好；未成功运行仍保留。}{fig:computation}
\FloatBarrier


\FloatBarrier
\section{实现参数与完整记录}\label{sec:parameters}
下面集中列出前述实验的固定设置；最终观测没有参与这些参数的选择。
主合成比较使用三十个训练种子，半径比较的各条件使用十个。
合成区间按相应实验中的训练种子进行 $2{,}000$ 次自助重采样，给出 $95\%$ 百分位区间；
每个种子的测试批次先取平均。
真实数据分割相互重叠，因此给出均值与分割标准差。
所有达到时间上限和失败的运行均保留在相应分母中，实际计算时间也全部计入。
\begin{table}[htbp]\centering\small\setlength{\tabcolsep}{3.5pt}
\caption{参考统计数据。选择与分类器拟合样本数均为每类数量；最终样本列注明单位。重复次数为训练种子数或重叠的真实数据分割数。}\label{tab:data}
\begin{tabularx}{\linewidth}{@{}p{25mm}rYrr@{}}
\toprule 数据 & $D/d$ & \shortstack{选择 / 拟合\\每类样本数} & \shortstack{最终\\观测} & 重复次数 \\
\midrule
合成数据 & 12/3 & 20/20 & 1000/类 & 30 \\
WDBC & 30/3 & 20/20 & 489 总数 & 20 \\
Ionosphere & 34/3 & 20/20 & 271 总数 & 20 \\
\bottomrule\end{tabularx}\end{table}
\FloatBarrier
\begin{table}[htbp]\centering\small\setlength{\tabcolsep}{3.5pt}
\caption{参考优化参数。每行分别在均值与混合备择上评价。}\label{tab:computation-settings}
\begin{tabularx}{\linewidth}{@{}rrrrY@{}}
\toprule $D$ & $d$ & 源点/类 & 种子数 & 备择 \\
\midrule
8 & 2 & 8 & 10 & 均值、混合 \\
12 & 3 & 8 & 10 & 均值、混合 \\
20 & 3 & 10 & 5 & 均值、混合 \\
40 & 3 & 10 & 5 & 均值、混合 \\
80 & 3 & 10 & 5 & 均值、混合 \\
\bottomrule\end{tabularx}\end{table}
\FloatBarrier
主实验两类半径均为 0.25。每个合成种子分别使用 20 个备择和 20 个零假设批次，每批每类 40 个观测。半径分析在每个机制和半径上使用 10 个种子。

中心化与样本标准差缩放只用合并的选择观测拟合，之后固定。共同 RBF 带宽使用其平方距离中位数。后续 SVM 使用 $C=1$，中位带宽只用自身拟合样本确定。鲁棒分类器使用与选择阶段相同的半合页损失和半径。

稀疏加性目标从带权核均值中减去 $0.1$ 倍方差估计，权重满足单位范数与稀疏约束。$D\le12$ 时枚举组合，更大问题使用前向选择与至多十次改进交换。实现采用论文中的负方差项，而不是核查脚本版本中的正项；测试带宽不使用测试标签重新拟合。该系数未进行交叉验证。L1 逻辑排序使用 liblinear、$C=1$ 与至多 $2{,}000$ 次迭代。贪心在时间上限时保留最后完成的阶段，并用 MMD 筛选补足余下变量。

统计实验每种方法的选择预算为 30 秒，半径为 0.25，数值间隙容差为 $10^{-6}+10^{-5}|\mathrm{UB}|$。较早的全局搜索以 MMD 筛选启动，并使用五轮分数起始求解。鲁棒随机化预测使用拟合的类别概率；SVM 的确定性判断以零为阈值，并列时判为类别零。

\begin{table}[htbp]\centering\small\setlength{\tabcolsep}{3.5pt}
\caption{单独的当前搜索确认实验的固定设置。}\label{tab:current-settings}
\begin{tabularx}{\linewidth}{@{}p{36mm}Y@{}}
\toprule 项目 & 设置 \\
\midrule
配置与数据 & 四份数据上的十六个配置；$(D,d)=(7,2),(10,3)$；每类六个观测；种子 19、23。 \\
损失与半径 & 半合页或逻辑损失；半径为 $(0.12,0.20)$ 或 $(0.30,0.45)$。 \\
生成坐标 & 类别零为高斯；类别一两个均值平移 1.4、0.9；噪声尺度 0.3--3；不做拟合标准化。 \\
共同起始组合 & $|\bar x_{0j}-\bar x_{1j}|/(s_{0j}+s_{1j}+10^{-12})$; 标准差用总体分母，并列按索引；不使用真实信号身份。 \\
时间与精度 & 60 秒，一个线程；$\varepsilon=2\times10^{-5}$，相对容差为零；内层/主问题目标为 $\varepsilon/4$。 \\
组合预测 & GP: 均值 $1/2$，协方差 $\frac14 e^{-\|z-z'\|_1/\max(1,d)}$，正则量 $\max\{10^{-8},(\varepsilon/(4B))^2\}$。 \\
候选与刷新 & 至多八个候选；均值减标准差评分；每完整求解三次或没有候选时刷新主问题；无分数根节点调用。 \\
\bottomrule\end{tabularx}\end{table}
\FloatBarrier
高斯过程（GP）预测归一化风险及其模型不确定性；归一化常数在半合页时为 $B=1$，逻辑损失时为 $B=2\log2$。确认实验的组合数使用求解器 visited 计数。受限的历史运行可能包含未达精度的末次评价，但当前确认的每次评价均达到 $\varepsilon/4$。实际精度重算与主问题时间保留在原始轨迹中。每配置计时三次，轮换执行顺序。

参考计算使用 AMD Ryzen 9 7945HX with Radeon Graphics、一个求解线程和 SciPy 内置的 HiGHS 1.12.0；逻辑损失核查使用 CVXPY 1.9.3 与 Clarabel 0.11.1。

较早研究的 1640 条记录包含 3 次失败和 189 次达到时间上限。完整结果、原尝试与重跑均保留在 \path{results/research_20261003} 中。辅助指标、所选组合和优化轨迹可在其中的 \path{per_trial.csv} 与 JSON 记录中查阅。开发对照、集成核验和较早的分式比较集中于 \path{docs/experiment_development_record.md}；当前确认单独保存在 \path{reports/verification/20261010/next_selection/current} 下。报告清单记录了实际使用的输入。


\begin{thebibliography}{9}
\bibitem{wang} J. Wang, S. S. Dey, and Y. Xie.
\emph{Variable Selection for Kernel Two-Sample Tests}. 2024.
\href{https://arxiv.org/abs/2302.07415}{arXiv:2302.07415v4}.
代码核查版本为
\href{https://github.com/WalterBabyRudin/MMDVar_Selection}{MMDVar\_Selection}
的提交 \texttt{031f740e9ca6}。
\bibitem{wdbc} UCI Machine Learning Repository.
\emph{Breast Cancer Wisconsin (Diagnostic)}. \href{https://doi.org/10.24432/C5DW2B}{DOI:10.24432/C5DW2B}.
\bibitem{ionosphere} UCI Machine Learning Repository.
\emph{Ionosphere}. \href{https://doi.org/10.24432/C5W01B}{DOI:10.24432/C5W01B}.
\end{thebibliography}
\end{document}
